% =============================================================================
% StructureTower 数論シリーズ NT2: Herbrand 関数と上付番号
% LuaLaTeX document — Lean 4 形式化の数学的解説
% =============================================================================
\documentclass[a4paper,11pt]{ltjsarticle}

% --- Core packages ---
\usepackage{luatexja-fontspec}
\usepackage{amsmath,amssymb,amsthm}
\usepackage{mathtools}
\usepackage{tikz}
\usepackage{tikz-cd}
\usetikzlibrary{cd,arrows,decorations.pathmorphing,matrix,positioning,shapes,calc}
\usepackage{enumitem}
\usepackage{hyperref}
\usepackage{cleveref}
\usepackage{tcolorbox}
\usepackage{listings}
\usepackage{xcolor}
\usepackage{geometry}
\usepackage{fancyhdr}
\usepackage{float}
\usepackage{booktabs}
\usepackage{array}

% --- Page geometry ---
\geometry{margin=2.5cm, top=3cm, bottom=3cm}

% --- Fonts ---
\setmainjfont{Noto Serif JP}
\setsansjfont{Noto Sans JP}

% --- Colors ---
\definecolor{leanblue}{HTML}{2563EB}
\definecolor{leanbg}{HTML}{F8FAFC}
\definecolor{leancomment}{HTML}{6B7280}
\definecolor{leankeyword}{HTML}{7C3AED}
\definecolor{leanstring}{HTML}{059669}
\definecolor{accentcolor}{HTML}{1E40AF}
\definecolor{theoremcolor}{HTML}{EFF6FF}
\definecolor{definitioncolor}{HTML}{F0FDF4}
\definecolor{remarkcolor}{HTML}{FFF7ED}

% --- Hyperref ---
\hypersetup{
  colorlinks=true,
  linkcolor=accentcolor,
  urlcolor=leanblue,
  citecolor=accentcolor,
  pdftitle={StructureTower 数論シリーズ NT2: Herbrand 関数と上付番号},
  pdfauthor={su},
}

% --- Theorem environments ---
\theoremstyle{definition}
\newtheorem{definition}{定義}[section]
\newtheorem{example}[definition]{例}

\theoremstyle{plain}
\newtheorem{theorem}[definition]{定理}
\newtheorem{lemma}[definition]{補題}
\newtheorem{proposition}[definition]{命題}
\newtheorem{corollary}[definition]{系}

\theoremstyle{remark}
\newtheorem{remark}[definition]{注意}

% --- Lean code environment ---
\lstdefinelanguage{lean4}{
  morekeywords={def,theorem,lemma,example,instance,class,structure,
    where,let,in,do,if,then,else,match,with,fun,return,
    import,open,namespace,section,end,variable,
    inductive,abbrev,noncomputable,private,protected,
    sorry,admit,have,show,suffices,calc,by,exact,
    apply,intro,intros,constructor,cases,induction,
    simp,rw,rfl,ext,ring,linarith,omega,norm_num,
    decide,trivial,assumption,contradiction,
    Type,Prop,Sort,Set,true,false,
    attribute,deriving},
  sensitive=true,
  morecomment=[l]{--},
  morecomment=[n]{/-}{-/},
  morestring=[b]",
}

\lstnewenvironment{leancode}[1][]{%
  \lstset{
    language=lean4,
    basicstyle=\ttfamily\small,
    keywordstyle=\color{leankeyword}\bfseries,
    commentstyle=\color{leancomment}\itshape,
    stringstyle=\color{leanstring},
    backgroundcolor=\color{leanbg},
    frame=single,
    rulecolor=\color{leanblue!30},
    framesep=8pt,
    xleftmargin=12pt,
    xrightmargin=12pt,
    breaklines=true,
    breakatwhitespace=true,
    showstringspaces=false,
    tabsize=2,
    captionpos=b,
    numbers=left,
    numberstyle=\tiny\color{leancomment},
    numbersep=8pt,
    aboveskip=1em,
    belowskip=1em,
    inputencoding=utf8,
    extendedchars=true,
    literate=
      {->}{{\(\to\)}}1
      {<->}{{\(\leftrightarrow\)}}1
      {<=}{{\(\leq\)}}1
      {>=}{{\(\geq\)}}1
      {/=}{{\(\neq\)}}1
      {:=}{{:=}}2,
    #1
  }%
}{}

% --- Tcolorbox styles ---
\tcbuselibrary{skins,breakable}

\newtcolorbox{proofstrategy}{
  colback=remarkcolor,
  colframe=orange!60!black,
  title={\textbf{証明戦略}},
  fonttitle=\sffamily,
  boxrule=0.5pt,
  arc=3pt,
}

\newtcolorbox{mathinsight}{
  colback=theoremcolor,
  colframe=accentcolor,
  title={\textbf{数学的洞察}},
  fonttitle=\sffamily,
  boxrule=0.5pt,
  arc=3pt,
}

\newtcolorbox{keyidea}{
  colback=definitioncolor,
  colframe=green!50!black,
  title={\textbf{核心的アイデア}},
  fonttitle=\sffamily,
  boxrule=0.5pt,
  arc=3pt,
}

% --- Math macros ---
\newcommand{\ST}{\mathsf{ST}}
\newcommand{\Hom}{\mathrm{Hom}}
\newcommand{\id}{\mathrm{id}}
\newcommand{\op}{\mathrm{op}}
\newcommand{\Nat}{\mathbb{N}}
\newcommand{\ZZ}{\mathbb{Z}}
\newcommand{\QQ}{\mathbb{Q}}
\newcommand{\Gal}{\mathrm{Gal}}
\newcommand{\Aut}{\mathrm{Aut}}
\newcommand{\ord}{\mathbb{N}^{\mathrm{op}}}
\newcommand{\iInf}{\displaystyle\bigcap}
\newcommand{\Phi}{\varPhi}
\newcommand{\varph}{\varphi}

% --- Header/Footer ---
\pagestyle{fancy}
\fancyhf{}
\renewcommand{\headrulewidth}{0.4pt}
\fancyhead[L]{\small\sffamily\nouppercase{\leftmark}}
\fancyhead[R]{\small\sffamily Lean 4 形式化 --- NT2}
\fancyfoot[C]{\thepage}

% =============================================================================
\begin{document}

% --- Title ---
\title{%
  {\LARGE\sffamily\bfseries StructureTower 数論シリーズ NT2}\\[0.5em]
  {\large\sffamily Herbrand 関数と上付番号：reindex の数論的具体化}\\[0.3em]
  {\normalsize\sffamily Lean 4 / Mathlib4 による形式化とその数学的解説}%
}
\author{su}
\date{\today}
\maketitle
\thispagestyle{fancy}

% --- AI Disclosure ---
\begin{center}
\small
\textit{AI assistance disclosure:}\\
Lean ソースコードは Claude (Anthropic) で骨格を生成し、
Codex (OpenAI) で修正した。\\
\TeX\ 文書は Claude Code (Anthropic) で生成した。\\
著者による加筆・修正は行っていない。\\
内容の正確性は保証されず、誤りがあれば著者の責任である。
\end{center}

\vspace{1em}

% --- Abstract ---
\begin{abstract}
本稿は Lean 4 / Mathlib4 による形式化ファイル
\texttt{StructureTower\_NT2\_HerbrandUpperNumbering.lean} の数学的解説である。
NT1 で構築した分岐塔 $\mathrm{ramificationTower}(\mathrm{rd}) \in
\textsc{StructureTower}\,\mathbb{N}^{\mathrm{op}}\,G$
（下付番号フィルトレーション $G_0 \supseteq G_1 \supseteq G_2 \supseteq \cdots$）に対し、
\textbf{Herbrand 関数} $\varPhi\colon \mathbb{N} \to \mathbb{N}$
による添字の変換（上付番号）を導入する。
上付番号は下付番号の致命的欠陥
—「部分群への制限と非整合」—を解消し、$H^u = H \cap G^u$ という基本性質を持つ。
\textsc{StructureTower} の視点では、上付番号塔は
$\mathrm{reindex}\,(\text{herbrandReindex})\,(\mathrm{ramificationTower})$
であり、L1 で定義した抽象的 \texttt{reindex} 操作の\textbf{最初の数論的インスタンス}となる。
さらに次数付き塔（Graded Tower）理論、ジャンプ集合の定義、
および Hasse-Arf 定理（アーベル拡大における上付番号ジャンプの整数性）の
定式化を行う。
\end{abstract}

\tableofcontents
\newpage

% =============================================================================
\section{はじめに：下付番号の欠陥と上付番号の動機}
% =============================================================================

\subsection{NT1 の振り返り}

NT1 では局所体 $L/K$，ガロア群 $G = \mathrm{Gal}(L/K)$，
均一化元 $\pi$，付値 $v_L$ に対する\textbf{下付番号分岐フィルトレーション}
\[
  G_i = \bigl\{\,\sigma \in G \;\big|\; v_L(\sigma(\pi) - \pi) \geq i + 1\,\bigr\}
\]
を $\textsc{StructureTower}\,\mathbb{N}^{\mathrm{op}}\,G$ として形式化した。
これは正規部分群の減少列
\[
  G = G_0 \supseteq G_1 \supseteq G_2 \supseteq \cdots \supseteq \{1\}
\]
を与え、添字 $i$ が「分岐の深さ」を直接表す。

\subsection{下付番号の致命的欠陥}

\begin{figure}[H]
\centering
\begin{tikzpicture}[>=stealth, node distance=1.6cm]
  % Left: G filtration
  \node[font=\bfseries] (Gtitle) {$G$ の下付番号};
  \node[below=0.4cm of Gtitle] (G0) {$G_0$};
  \node[below of=G0] (G1) {$G_1$};
  \node[below of=G1] (G2) {$G_2$};
  \node[below of=G2] (Gdots) {$\vdots$};
  \draw[right hook->] (G1) -- (G0) node[midway,right] {\small $\trianglelefteq$};
  \draw[right hook->] (G2) -- (G1) node[midway,right] {\small $\trianglelefteq$};
  \draw[right hook->] (Gdots) -- (G2);

  % Middle: H restriction
  \node[right=3cm of Gtitle, font=\bfseries] (Htitle) {$H \leq G$ の制限};
  \node[below=0.4cm of Htitle] (H0) {$H \cap G_0$};
  \node[below of=H0] (H1) {$H \cap G_1$};
  \node[below of=H1] (H2) {$H \cap G_2$};
  \node[below of=H2] (Hdots) {$\vdots$};
  \draw[right hook->] (H1) -- (H0);
  \draw[right hook->] (H2) -- (H1);
  \draw[right hook->] (Hdots) -- (H2);

  % Right: H's own lower numbering
  \node[right=3cm of Htitle, font=\bfseries] (Howtitle) {$H$ 自身の下付番号};
  \node[below=0.4cm of Howtitle] (H0own) {$H_0$};
  \node[below of=H0own] (H1own) {$H_1$};
  \node[below of=H1own] (H2own) {$H_2$};
  \node[below of=H2own] (Hdotsown) {$\vdots$};
  \draw[right hook->] (H1own) -- (H0own);
  \draw[right hook->] (H2own) -- (H1own);
  \draw[right hook->] (Hdotsown) -- (H2own);

  % Inequality signs
  \node[right=0.8cm of H0, font=\Large, color=red] {$\neq$};
  \node[right=0.8cm of H1, font=\Large, color=red] {$\neq$};
  \node[right=0.8cm of H2, font=\Large, color=red] {$\neq$};
\end{tikzpicture}
\caption{下付番号の欠陥：$H \cap G_i$ は一般に $H$ 自身の分岐フィルトレーション $H_i$ と一致しない。
添字のずれが生じる。}
\label{fig:lower-numbering-defect}
\end{figure}

下付番号の問題点：$H \leq G$ に対し、$H$ を $G$ の分岐フィルトレーションで制限した
$H \cap G_i$ は、一般に $H$ 自身の分岐指標で決まる分岐フィルトレーション $H_i$
と一致しない（図~\ref{fig:lower-numbering-defect}）。

\subsection{解決策：Herbrand 関数}

Herbrand の解決策は添字を変換することである。
\textbf{Herbrand 関数}（連続版）：
\[
  \varphi(t) = \int_0^t \frac{1}{[G_0 : G_s]}\,ds
\]
（$G_s$ の指数の逆数を積分する）により得られる\textbf{上付番号}は：
\[
  G^u = G_{\psi(u)}, \quad \psi = \varphi^{-1}
\]
と定義され、\textbf{制限と可換}になる：
\[
  H^u = H \cap G^u \quad (\text{上付番号は制限と可換！})
\]

\begin{keyidea}
\textsc{StructureTower} の視点での核心：
\[
  \text{上付番号塔} = \mathrm{reindex}\,\varPhi\,h_{\varPhi}\,(\mathrm{ramificationTower}\,\mathrm{rd})
\]
Herbrand 関数 $\varPhi$ は L1 で定義した抽象的 \texttt{reindex} の
\textbf{具体的数論的インスタンス}である。
「reindex の選び方に算術的内容がある」ことを示す最初の例。
\end{keyidea}

\subsection{NT2 の構成}

\begin{center}
\begin{tikzpicture}
  \node[draw,rounded corners,fill=blue!8,text width=3.5cm,align=center] (nt21)
    {\S NT2-1\\次数付き塔};
  \node[draw,rounded corners,fill=green!8,text width=3.5cm,align=center,
    right=1cm of nt21] (nt22)
    {\S NT2-2\\Herbrand 関数};
  \node[draw,rounded corners,fill=orange!8,text width=3.5cm,align=center,
    right=1cm of nt22] (nt23)
    {\S NT2-3\\上付番号塔};
  \node[draw,rounded corners,fill=purple!8,text width=3.5cm,align=center,
    below=1cm of nt22] (nt24)
    {\S NT2-4\\商との整合性};
  \node[draw,rounded corners,fill=red!8,text width=3.5cm,align=center,
    right=1cm of nt24] (nt25)
    {\S NT2-5\\Hasse-Arf 定式化};
  \draw[->] (nt21) -- (nt22);
  \draw[->] (nt22) -- (nt23);
  \draw[->] (nt23) -- (nt24);
  \draw[->] (nt22) -- (nt24);
  \draw[->] (nt24) -- (nt25);
\end{tikzpicture}
\end{center}

% =============================================================================
\section{\S NT2-1. 次数付き塔（Graded Tower）}
% =============================================================================

\subsection{連続レベル間の差分}

\textsc{StructureTower} は「各レベル以下の元の集合」を与えるが、
「ちょうどそのレベルにある元」の集合を直接記述する概念がなかった。
次数付き塔理論はこの空白を埋める。

\begin{definition}[次数片（gradedPiece）]\label{def:graded-piece}
  $\mathbb{N}^{\mathrm{op}}$ 添字の塔 $T \in \textsc{StructureTower}\,\mathbb{N}^{\mathrm{op}}\,\alpha$
  に対し、\textbf{次数 $n$ の次数片}を
  \[
    \mathrm{gradedPiece}(T,n)
    := \mathrm{level}(\tilde{n}) \setminus \mathrm{level}(\widetilde{n+1})
    \quad (n \in \mathbb{N})
  \]
  で定める（$\tilde{n} = \mathrm{toDual}(n) \in \mathbb{N}^{\mathrm{op}}$）。

  分岐塔では $\mathrm{gradedPiece}(\mathrm{ramificationTower}\,\mathrm{rd},\,n) = G_n \setminus G_{n+1}$
  —「分岐の深さがちょうど $n$ の自己同型」の集合。
\end{definition}

\begin{figure}[H]
\centering
\begin{tikzpicture}[>=stealth]
  % Number line for levels
  \draw[->] (0,0) -- (9,0) node[right] {$n$ (深さ)};
  % Bars for levels
  \foreach \i/\col in {0/blue!40, 1/blue!30, 2/blue!20, 3/blue!10} {
    \fill[\col] (\i*2, 0.1) rectangle (\i*2+1.8, 1.5-\i*0.3);
  }
  \draw[thick] (0,0.1) rectangle (1.8,1.5);
  \draw[thick] (2,0.1) rectangle (3.8,1.2);
  \draw[thick] (4,0.1) rectangle (5.8,0.9);
  \draw[thick] (6,0.1) rectangle (7.8,0.6);
  % Labels
  \node at (0.9, 0.75) {\small $G_0$};
  \node at (2.9, 0.65) {\small $G_1$};
  \node at (4.9, 0.5) {\small $G_2$};
  \node at (6.9, 0.35) {\small $G_3$};
  % Graded pieces
  \node[above, color=red] at (0.9, 1.5) {\small $\mathrm{grd}(0)$};
  \node[above, color=red] at (2.9, 1.2) {\small $\mathrm{grd}(1)$};
  % Arrows showing difference
  \draw[<->, red, thick] (1.9, 1.5) -- (1.9, 1.2) node[midway, right] {\small $G_0\setminus G_1$};
  \draw[<->, red, thick] (3.9, 1.2) -- (3.9, 0.9) node[midway, right] {\small $G_1\setminus G_2$};
\end{tikzpicture}
\caption{次数片 $\mathrm{gradedPiece}(T,n) = G_n \setminus G_{n+1}$ の図示。
各バーの高さが部分群の「大きさ」を表す。差分が次数片。}
\label{fig:graded-piece}
\end{figure}

\subsection{次数片の基本性質}

\begin{proposition}[成員条件]\label{prop:mem-graded-piece}
  \[
    x \in \mathrm{gradedPiece}(T,n) \iff
    x \in T.\mathrm{level}(\tilde{n}) \;\land\; x \notin T.\mathrm{level}(\widetilde{n+1})
  \]
\end{proposition}

\begin{proof}
  差分集合の定義（\texttt{Set.mem\_diff}）から直ちに従う。
\end{proof}

\begin{theorem}[次数片の互いに素性]\label{thm:graded-piece-disjoint}
  $m \neq n$ ならば $\mathrm{gradedPiece}(T,m) \cap \mathrm{gradedPiece}(T,n) = \emptyset$。
\end{theorem}

\begin{proof}
  $m < n$ の場合を示す（$m > n$ は対称的）。
  $x \in \mathrm{gradedPiece}(T,m)$ ならば $x \notin T.\mathrm{level}(\widetilde{m+1})$。
  一方 $n \geq m+1$ より、単調性から
  $T.\mathrm{level}(\tilde{n}) \subseteq T.\mathrm{level}(\widetilde{m+1})$
  （$\mathbb{N}^{\mathrm{op}}$ では $\tilde{n} \leq \widetilde{m+1}$）。
  したがって $x \in \mathrm{gradedPiece}(T,n)$ ならば $x \in T.\mathrm{level}(\tilde{n})
  \subseteq T.\mathrm{level}(\widetilde{m+1})$ となり矛盾。
\end{proof}

\begin{figure}[H]
\centering
\begin{tikzpicture}[>=stealth]
  \draw[->] (-0.5, 0) -- (8.5, 0) node[right] {$n$};
  % gradedPiece intervals (disjoint)
  \fill[blue!30] (0,0.2) rectangle (1.5, 0.8);
  \fill[green!30] (2,0.2) rectangle (3.5, 0.8);
  \fill[orange!30] (4,0.2) rectangle (5.5, 0.8);
  \fill[red!30] (6,0.2) rectangle (7.5, 0.8);
  \node at (0.75, 0.5) {\small $\mathrm{grd}(0)$};
  \node at (2.75, 0.5) {\small $\mathrm{grd}(1)$};
  \node at (4.75, 0.5) {\small $\mathrm{grd}(2)$};
  \node at (6.75, 0.5) {\small $\mathrm{grd}(3)$};
  % Disjoint labels
  \foreach \x in {1.75, 3.75, 5.75} {
    \node at (\x, 0.5) {\Large $\cap$};
    \node at (\x, 0.2) {\small $=\emptyset$};
  }
  \node[above] at (4, 0.8) {\bfseries 互いに素};
\end{tikzpicture}
\caption{定理~\ref{thm:graded-piece-disjoint}：次数片は互いに素。
異なる次数の元は同時に2つの次数片に属せない。}
\label{fig:graded-disjoint}
\end{figure}

\begin{proposition}[和集合との関係]\label{prop:union-graded}
  \[
    \bigcup_{n \in \mathbb{N}} \mathrm{gradedPiece}(T,n) \subseteq T.\mathrm{union}
  \]
\end{proposition}

\begin{proof}
  $x \in \mathrm{gradedPiece}(T,n)$ ならば $x \in T.\mathrm{level}(\tilde{n})$。
  したがって $x \in \bigcup_i T.\mathrm{level}(i) = T.\mathrm{union}$。
\end{proof}

\begin{mathinsight}
命題~\ref{prop:union-graded} の逆は一般に成立しない。
$T.\mathrm{union}$ の元 $x$ が $T.\mathrm{level}(\tilde{n})$ に属しても、
$T.\mathrm{level}(\widetilde{n+1})$ からも「脱出しない」場合、
$x$ はどの次数片にも属さない（「最大次数」を持たない元）。

分離条件（$\bigcap_n G_n = \{1\}$）のもとでは、
$1$ を除く各元は「最後に属するレベル」が存在し、
$\bigcup_n \mathrm{gradedPiece}(T,n) = T.\mathrm{union} \setminus \{1\}$
という完全分解が成立する（有限群の場合）。
\end{mathinsight}

\begin{proposition}[分岐塔での具体化]\label{prop:graded-ramification}
  \[
    \mathrm{gradedPiece}(\mathrm{ramificationTower}(\mathrm{rd}),\,n)
    = \lvert G_n \rvert \setminus \lvert G_{n+1} \rvert
  \]
  （$G_n = \mathrm{rd.groups}(n)$、台集合の差分）。
\end{proposition}

\begin{proof}
  \texttt{ramificationTower\_level} を展開すれば定義から直接従う。
\end{proof}

\begin{mathinsight}
次数付き塔理論が \textsc{OrderHom} を超える3点：
\begin{enumerate}
  \item \textbf{隣接レベルの概念}：$\mathbb{N}^{\mathrm{op}}$ の離散構造（後者関数 $n \mapsto n+1$）
    に依存する。一般の前順序では「隣接する」という概念がない。
  \item \textbf{直和分解}：$T.\mathrm{union} = \bigsqcup_n \mathrm{gradedPiece}(T,n)$
    （適切な条件下）は非自明な定理。
  \item \textbf{分岐指標}：$b(\sigma) = $ 「$\sigma$ が属する次数片のインデックス」
    として解釈できる。
\end{enumerate}
\end{mathinsight}

% =============================================================================
\section{\S NT2-2. Herbrand 関数}
% =============================================================================

\subsection{連続版から離散版へ}

連続版の Herbrand 関数
\[
  \varphi(t) = \int_0^t [G_0 : G_s]^{-1}\,ds
\]
を $\mathbb{N}$ 上で離散化する。
$s \in [i, i+1)$ の範囲では $G_s = G_i$（一定）なので：
\[
  \varphi(n) = \sum_{i=1}^{n} \frac{|G_i|}{|G_0|}
\]

Lean での扱いやすさのため，\textbf{スケーリング版}（$|G_0|$ 倍）を採用する：
\[
  \Phi(0) = 0, \qquad \Phi(n) = \sum_{i=1}^{n} |G_i| \quad (n \geq 1)
\]

\begin{figure}[H]
\centering
\begin{tikzpicture}[>=stealth, scale=0.85]
  \draw[->] (0,0) -- (7,0) node[right] {$n$ (下付番号)};
  \draw[->] (0,0) -- (0,5) node[above] {$\Phi(n)$ (スケール上付)};
  % Example: |G_1|=4, |G_2|=2, |G_3|=1, ...
  \foreach \x/\y in {0/0, 1/4, 2/6, 3/7} {
    \fill[blue] (\x, \y*0.6) circle (3pt);
  }
  \draw[blue, thick] (0,0) -- (1, 2.4) -- (2, 3.6) -- (3, 4.2);
  \node[right, blue] at (3, 4.2) {\small $\Phi(n) = \Sigma_{i=1}^n |G_i|$};
  % Slope indicators
  \node[above, color=red] at (0.5, 1.2) {\small $|G_1|=4$};
  \node[above, color=red] at (1.5, 3.0) {\small $|G_2|=2$};
  \node[above, color=red] at (2.5, 3.9) {\small $|G_3|=1$};
  % Reference line (identity)
  \draw[dashed, gray] (0,0) -- (3, 1.8) node[right] {\small $\Phi(n) = n$};
\end{tikzpicture}
\caption{スケーリング Herbrand 関数 $\Phi$ の例（$|G_1|=4$，$|G_2|=2$，$|G_3|=1$）。
$\Phi$ は単調増加し、$\Phi(n) \geq n$（各 $|G_i| \geq 1$ のとき）。
実際の Herbrand 関数は $\varphi(n) = \Phi(n)/|G_0|$。}
\label{fig:herbrand-function}
\end{figure}

\subsection{scaledHerbrand の定義と基本性質}

\begin{definition}[スケーリング Herbrand 関数]\label{def:scaled-herbrand}
  $\mathrm{groupSize}\colon \mathbb{N} \to \mathbb{N}$（各レベルの群のサイズを抽象化）に対し、
  \[
    \mathrm{scaledHerbrand}(\mathrm{gs})\colon \mathbb{N} \to \mathbb{N}
  \]
  を漸化式
  \[
    \Phi(0) = 0, \qquad \Phi(n+1) = \Phi(n) + \mathrm{gs}(n+1)
  \]
  で定める。\textbf{Finset.sum による表現}：
  \[
    \Phi(n) = \sum_{i \in \{0,\ldots,n-1\}} \mathrm{gs}(i+1)
    = \mathrm{gs}(1) + \mathrm{gs}(2) + \cdots + \mathrm{gs}(n)
  \]
\end{definition}

\begin{leancode}
def scaledHerbrand (groupSize : N -> N) : N -> N
  | 0     => 0
  | n + 1 => scaledHerbrand groupSize n + groupSize (n + 1)

-- Finset.sum による明示的表現
theorem scaledHerbrand_eq_sum (gs : N -> N) (n : N) :
    scaledHerbrand gs n = (Finset.range n).sum (fun i => gs (i + 1)) := by
  induction n with
  | zero     => simp
  | succ k ih => simp [Finset.sum_range_succ, ih]
\end{leancode}

\begin{theorem}[単調性]\label{thm:scaled-herbrand-mono}
  $\mathrm{scaledHerbrand}(\mathrm{gs})$ は単調：$m \leq n \Rightarrow \Phi(m) \leq \Phi(n)$。
\end{theorem}

\begin{proof}
  後者に関して $\Phi(n+1) = \Phi(n) + \mathrm{gs}(n+1) \geq \Phi(n)$（$\mathbb{N}$ での等式）。
  \texttt{monotone\_nat\_of\_le\_succ} で一般の $m \leq n$ に帰着。
\end{proof}

\begin{theorem}[狭義単調性]\label{thm:scaled-herbrand-strict-mono}
  $\mathrm{gs}(i) > 0$（$i > 0$ のとき）ならば $\Phi$ は狭義単調：
  $m < n \Rightarrow \Phi(m) < \Phi(n)$。

  数論的意味：各 $G_i \neq \emptyset$（部分群は必ず単位元を含む）より $|G_i| \geq 1$。
  よって実際の分岐データでは Herbrand 関数は狭義単調 $\Rightarrow$ 単射
  $\Rightarrow$ 逆関数 $\psi = \Phi^{-1}$ が存在する。
\end{theorem}

\begin{proof}
  \[
    \Phi(n) < \Phi(n) + \mathrm{gs}(n+1) = \Phi(n+1)
  \]
  （$\mathrm{gs}(n+1) > 0$ より）。\texttt{strictMono\_nat\_of\_lt\_succ} で帰着。
\end{proof}

\subsection{herbrandReindex: ℕᵒᵈ → ℕᵒᵈ への変換}

\begin{figure}[H]
\centering
\begin{tikzcd}[column sep=3cm, row sep=2cm]
  \mathbb{N}
    \arrow[r, "\Phi", "\text{単調}"']
  & \mathbb{N} \\
  \mathbb{N}^{\mathrm{op}}
    \arrow[u, "\mathrm{ofDual}"]
    \arrow[r, "\text{herbrandReindex}"', dashrightarrow]
  & \mathbb{N}^{\mathrm{op}}
    \arrow[u, "\mathrm{ofDual}"']
\end{tikzcd}
\caption{\texttt{herbrandReindex} の構成。
$\mathbb{N}$ 上の単調な $\Phi$ から $\mathbb{N}^{\mathrm{op}}$ 上の単調写像を誘導する。
$\mathbb{N}^{\mathrm{op}}$ では大小関係が逆転するので、$\Phi$ が単調 $\Rightarrow$
\texttt{herbrandReindex} も単調。}
\label{fig:herbrand-reindex}
\end{figure}

\begin{definition}[herbrandReindex]\label{def:herbrand-reindex}
  \[
    \mathrm{herbrandReindex}(\mathrm{gs})\colon \mathbb{N}^{\mathrm{op}} \to \mathbb{N}^{\mathrm{op}},
    \qquad n \mapsto \mathrm{toDual}(\Phi(\mathrm{ofDual}(n)))
  \]
\end{definition}

\begin{proposition}\label{prop:herbrand-reindex-mono}
  \texttt{herbrandReindex gs} は単調写像（$\mathbb{N}^{\mathrm{op}} \to \mathbb{N}^{\mathrm{op}}$）。
\end{proposition}

\begin{proof}
  $i \leq j$ in $\mathbb{N}^{\mathrm{op}}$ $\Leftrightarrow$ $\mathrm{ofDual}(j) \leq \mathrm{ofDual}(i)$ in $\mathbb{N}$
  $\Rightarrow$ $\Phi(\mathrm{ofDual}(j)) \leq \Phi(\mathrm{ofDual}(i))$ （$\Phi$ の単調性）
  $\Leftrightarrow$ $\mathrm{toDual}(\Phi(\mathrm{ofDual}(i))) \leq \mathrm{toDual}(\Phi(\mathrm{ofDual}(j)))$
  in $\mathbb{N}^{\mathrm{op}}$。
\end{proof}

% =============================================================================
\section{\S NT2-3. 上付番号塔}
% =============================================================================

\subsection{reindex による構成}

\begin{definition}[上付番号塔]\label{def:upper-numbering-tower}
  分岐データ $\mathrm{rd}$，サイズ関数 $\mathrm{gs}$ に対し、
  \textbf{上付番号塔}を
  \begin{align*}
    \mathrm{upperNumberingTower}(\mathrm{rd},\mathrm{gs})
    &:= \mathrm{reindex}\,(\mathrm{herbrandReindex}\,\mathrm{gs})\,
       h_\Phi\,(\mathrm{ramificationTower}\,\mathrm{rd}) \\
    &\in \textsc{StructureTower}\,\mathbb{N}^{\mathrm{op}}\,G
  \end{align*}
  で定める（$h_\Phi$：命題~\ref{prop:herbrand-reindex-mono} の単調性証明）。
\end{definition}

\begin{figure}[H]
\centering
\begin{tikzcd}[column sep=2.5cm, row sep=2cm]
  \mathbb{N}^{\mathrm{op}}
    \arrow[r, "\text{herbrandReindex}"]
    \arrow[rr, bend left=20, "\text{level (上付番号)}"]
  & \mathbb{N}^{\mathrm{op}}
    \arrow[r, "\text{level (下付番号)}"]
  & \mathcal{P}(G) \\
  u
    \arrow[r, mapsto]
  & \Phi(u)
    \arrow[r, mapsto]
  & G_{\Phi(u)}
\end{tikzcd}
\caption{上付番号塔の構成。上付番号 $u$ における level は
$G^u = G_{\Phi(u)}$（$\Phi$ を経由して下付番号の塔を参照する）。}
\label{fig:upper-numbering-tower}
\end{figure}

\begin{proposition}[level の明示的表示]\label{prop:upper-numbering-level}
  \[
    (\mathrm{upperNumberingTower}(\mathrm{rd},\mathrm{gs})).\mathrm{level}(n)
    = \lvert \mathrm{rd.groups}(\Phi(\mathrm{ofDual}(n)))\rvert
  \]
\end{proposition}

\begin{proof}
  \texttt{reindex\_level} と \texttt{ramificationTower\_level} を合成。
  $\mathrm{simp}$ で直接確認できる。
\end{proof}

\begin{proposition}[レベル 0 の整合性]\label{prop:upper-zero}
  \[
    (\mathrm{upperNumberingTower}(\mathrm{rd},\mathrm{gs})).\mathrm{level}(\tilde{0})
    = \lvert \mathrm{rd.groups}(0)\rvert = \lvert G_0\rvert
  \]
\end{proposition}

\begin{proof}
  $\Phi(0) = 0$（\texttt{scaledHerbrand\_zero}）より直ちに。
\end{proof}

\subsection{下付番号との比較}

\begin{theorem}[NatInclusion 関係]\label{thm:upper-nat-inclusion}
  $\mathrm{gs}(n) \geq 1$ ならば $\Phi(n) \geq n$ であり、
  \[
    \mathrm{NatInclusion}(\mathrm{upperNumberingTower}\,\mathrm{rd}\,\mathrm{gs},\;\mathrm{ramificationTower}\,\mathrm{rd})
  \]
  すなわち各 $n \in \mathbb{N}^{\mathrm{op}}$ で
  $G^n = G_{\Phi(n)} \subseteq G_n$（上付番号の方が細かいフィルタリング）。
\end{theorem}

\begin{proof}
  $\Phi(n) \geq n$ より $G_{\Phi(n)} \leq G_n$（$\mathrm{rd.antitone}$）。
  台集合の包含が従う。
\end{proof}

\begin{figure}[H]
\centering
\begin{tikzpicture}[>=stealth, scale=0.9]
  \draw[->] (0,0) -- (7,0) node[right] {$n$};
  % Lower numbering levels (wider)
  \foreach \i in {0,...,5} {
    \fill[blue!20] (\i*1.1, 0.1) rectangle (\i*1.1+1.0, 1.5-\i*0.2);
  }
  % Upper numbering levels (narrower -- more filtered)
  \foreach \i in {0,...,5} {
    \fill[red!30] (\i*1.1+0.1, 0.1) rectangle (\i*1.1+0.9, 1.3-\i*0.25);
    \draw[red, thick] (\i*1.1+0.1, 0.1) rectangle (\i*1.1+0.9, 1.3-\i*0.25);
  }
  % Legend
  \fill[blue!20] (0, -0.6) rectangle (0.5, -0.3);
  \node[right] at (0.6, -0.45) {\small 下付番号 $G_n$};
  \fill[red!30] (3, -0.6) rectangle (3.5, -0.3);
  \draw[red, thick] (3, -0.6) rectangle (3.5, -0.3);
  \node[right] at (3.6, -0.45) {\small 上付番号 $G^n \subseteq G_n$};
\end{tikzpicture}
\caption{定理~\ref{thm:upper-nat-inclusion}：上付番号の方が「細かい」。
各 $n$ で $G^n \subseteq G_n$。}
\label{fig:upper-vs-lower}
\end{figure}

\subsection{二重 reindex の合成則}

\begin{theorem}[二重 reindex]\label{thm:double-reindex}
  2つの Herbrand 関数を合成することに対応する塔の等式：
  \begin{align*}
    &\mathrm{reindex}\,(\mathrm{herbrandReindex}\,\mathrm{gs}_1)\,h_1\,
     (\mathrm{upperNumberingTower}\,\mathrm{rd}\,\mathrm{gs}_2) \\
    &= \mathrm{reindex}\,(\mathrm{herbrandReindex}\,\mathrm{gs}_2 \circ \mathrm{herbrandReindex}\,\mathrm{gs}_1)\,
       (h_2 \circ h_1)\,(\mathrm{ramificationTower}\,\mathrm{rd})
  \end{align*}
\end{theorem}

\begin{proof}
  \texttt{reindex\_comp} の直接的インスタンス。
  \texttt{simp} と \texttt{reindex} の定義展開で確認できる。
\end{proof}

\begin{mathinsight}
定理~\ref{thm:double-reindex} の意味：
「下付番号 $\to$ 上付番号(gs₂) $\to$ 再添字(gs₁)」という2段階の操作は、
「下付番号 $\to$ 合成 Herbrand $\mathrm{gs}_2 \circ \mathrm{gs}_1$ による添字変換」に等しい。
これは \textsc{StructureTower} の \texttt{reindex\_comp} が数論の文脈で
自動的に成立することを示す — 抽象的な関手の等式が算術的内容を自動的に含む。
\end{mathinsight}

% =============================================================================
\section{\S NT2-4. 商との整合性}
% =============================================================================

\subsection{部分群への制限}

\begin{definition}[部分群への制限塔]\label{def:restrict-subgroup}
  $T \in \textsc{StructureTower}\,\mathbb{N}^{\mathrm{op}}\,G$，$H \leq G$ に対し、
  \[
    \mathrm{restrictToSubgroup}(T,H).\mathrm{level}(n) := T.\mathrm{level}(n) \cap \lvert H\rvert
  \]
  は $\textsc{StructureTower}\,\mathbb{N}^{\mathrm{op}}\,G$ を定める。
\end{definition}

\begin{proposition}[制限は NatInclusion で以下]\label{prop:restrict-nat-inclusion}
  \[
    \mathrm{NatInclusion}(\mathrm{restrictToSubgroup}(T,H),\;T)
  \]
\end{proposition}

\begin{proof}
  $(T.\mathrm{level}(n) \cap \lvert H\rvert) \subseteq T.\mathrm{level}(n)$
  （\texttt{Set.inter\_subset\_left}）。
\end{proof}

\subsection{上付番号の商整合性}

下付番号が制限と非整合（図~\ref{fig:lower-numbering-defect}）なのに対し、
上付番号は整合する。これを定式化する。

\begin{definition}[上付番号整合条件]\label{def:upper-numbering-compatible}
  分岐データ $\mathrm{rd}_G$（$G$ 用），$\mathrm{rd}_H$（$H \leq G$ 用），
  サイズ関数 $\mathrm{gs}_G$，$\mathrm{gs}_H$，$H \leq G$ に対し、
  \textbf{上付番号整合条件}（\texttt{UpperNumberingCompatible}）を
  \[
    \forall n \in \mathbb{N}^{\mathrm{op}},\;
    (\mathrm{upperTower}(\mathrm{rd}_G,\mathrm{gs}_G)).\mathrm{level}(n) \cap \lvert H\rvert
    = (\mathrm{upperTower}(\mathrm{rd}_H,\mathrm{gs}_H)).\mathrm{level}(n)
  \]
  で定める。
\end{definition}

\begin{figure}[H]
\centering
\begin{tikzcd}[column sep=3cm, row sep=2.5cm]
  G^u = G_{\Phi_G(u)}
    \arrow[r, hookrightarrow]
    \arrow[d, "\cap\, H"']
  & G_{\Phi_G(u-1)}
    \arrow[d, "\cap\, H"] \\
  G^u \cap H
    \arrow[r, equal, "\text{整合条件}"]
  & H^u = H_{\Phi_H(u)}
\end{tikzcd}
\caption{上付番号の整合条件：$G$ の上付番号塔を $H$ に制限すると、
$H$ 自身の上付番号塔（適切な $\Phi_H$ で定義）と一致する。
これが Herbrand 変換の本質的性質。}
\label{fig:upper-compatibility}
\end{figure}

\begin{remark}
  定義~\ref{def:upper-numbering-compatible} の条件が成立するとき、
  上付番号は「正しい reindex」であると言う。
  下付番号（$\Phi = \mathrm{id}$ に対応する恒等 reindex）では
  一般にこの条件が失敗する。
\end{remark}

\subsection{下付番号は恒等 reindex}

\begin{theorem}[下付番号 = 恒等 reindex]\label{thm:lower-is-id-reindex}
  \[
    \mathrm{reindex}\,(\mathrm{id}\colon \mathbb{N}^{\mathrm{op}} \to \mathbb{N}^{\mathrm{op}})\,
    \mathrm{monotone\_id}\,(\mathrm{ramificationTower}\,\mathrm{rd})
    = \mathrm{ramificationTower}\,\mathrm{rd}
  \]
\end{theorem}

\begin{proof}
  \texttt{reindex\_id} の直接適用。
\end{proof}

\begin{mathinsight}
定理~\ref{thm:lower-is-id-reindex} は「下付番号 = $\Phi = \mathrm{id}$ による reindex」
であることを明示する。上付番号との対比：

\begin{center}
\renewcommand{\arraystretch}{1.3}
\begin{tabular}{lll}
\toprule
 & \textbf{reindex 関数} & \textbf{制限との整合} \\
\midrule
下付番号 & $\mathrm{id}$（恒等写像） & 一般に$\times$ \\
上付番号 & Herbrand $\Phi$（非自明） & $\bigcirc$（Hasse-Arf 条件下） \\
\bottomrule
\end{tabular}
\end{center}

\textsc{StructureTower} の言語では、「どの reindex を選ぶか」が
「部分構造との整合性」を決定する。Herbrand 関数は
この問いへの数論的な答え。
\end{mathinsight}

% =============================================================================
\section{\S NT2-5. Hasse-Arf 定理の定式化}
% =============================================================================

\subsection{Hasse-Arf 定理の概要}

\begin{figure}[H]
\centering
\begin{tikzpicture}[>=stealth, node distance=1.5cm]
  \node[draw, rounded corners, fill=blue!10, text width=4cm, align=center] (ha)
    {\textbf{Hasse-Arf 定理}\\(1930/1951)};
  \node[draw, rounded corners, fill=green!10, text width=5cm, align=center,
    right=1.5cm of ha] (stmt)
    {$L/K$ アーベルガロア拡大\\$\Rightarrow$\\上付番号のジャンプは$\in \mathbb{Z}$};
  \node[draw, rounded corners, fill=orange!10, text width=5cm, align=center,
    below=1.5cm of ha] (struct)
    {\textsc{StructureTower} 版\\$|G_0| \mid \Phi(b)$\\（$b$: 下付番号のジャンプ）};
  \draw[->] (ha) -- (stmt);
  \draw[->] (ha) -- (struct);
  \node[above=0.3cm of ha] {\small 類体論の深い定理};
\end{tikzpicture}
\caption{Hasse-Arf 定理の概要と \textsc{StructureTower} での定式化。}
\label{fig:hasse-arf-overview}
\end{figure}

下付番号のジャンプ（break）$b$ は定義から $\mathbb{N}$ の元（整数）だが、
上付番号のジャンプ $\varphi(b) = \Phi(b)/|G_0|$ は一般には有理数になる。
Hasse-Arf 定理はアーベル拡大でこれが整数になることを主張する。

\subsection{ジャンプ集合}

\begin{definition}[ジャンプ集合]\label{def:jump-set}
  $T \in \textsc{StructureTower}\,\mathbb{N}^{\mathrm{op}}\,\alpha$ に対し、
  \textbf{ジャンプ集合}を
  \[
    \mathrm{jumpSet}(T) := \bigl\{\,n \in \mathbb{N} \;\big|\;
    \mathrm{gradedPiece}(T,n) \neq \emptyset\,\bigr\}
  \]
  で定める。すなわち「段差のある」インデックスの集合。
\end{definition}

\begin{figure}[H]
\centering
\begin{tikzpicture}[>=stealth]
  \draw[->] (-0.5,0) -- (9,0) node[right] {$n$};
  % G_n steps
  \fill[blue!20] (0, 0.1) rectangle (2.5, 1.8);
  \fill[blue!20] (2.5, 0.1) rectangle (4.5, 1.2);
  \fill[blue!20] (4.5, 0.1) rectangle (6.5, 0.7);
  \fill[blue!20] (6.5, 0.1) rectangle (8.5, 0.3);
  \draw[thick] (0, 0.1) -- (0, 1.8) -- (2.5, 1.8);
  \draw[thick] (2.5, 1.2) -- (4.5, 1.2);
  \draw[thick] (4.5, 0.7) -- (6.5, 0.7);
  \draw[thick] (6.5, 0.3) -- (8.5, 0.3);
  % Jump indicators
  \draw[very thick, red] (2.5, 1.2) -- (2.5, 1.8);
  \draw[very thick, red] (4.5, 0.7) -- (4.5, 1.2);
  \draw[very thick, red] (6.5, 0.3) -- (6.5, 0.7);
  % Jump set markers
  \foreach \x/\lab in {2.5/b_1, 4.5/b_2, 6.5/b_3} {
    \fill[red] (\x, -0.15) circle (3pt);
    \node[below, red] at (\x, -0.15) {\small $\lab$};
  }
  \node[above] at (4.5, 1.8) {\small ジャンプ集合 $= \{b_1, b_2, b_3\}$};
  \draw[decorate, decoration={brace, amplitude=5pt}, red]
    (0, 2.0) -- (8.5, 2.0) node[midway, above=5pt] {\small $\mathrm{jumpSet}(T)$};
\end{tikzpicture}
\caption{ジャンプ集合の図示。赤い縦線の位置（$b_1, b_2, b_3$）が
$\mathrm{gradedPiece}$ が非空になるインデックス。}
\label{fig:jump-set}
\end{figure}

\begin{proposition}[成員条件]\label{prop:mem-jump-set}
  \[
    n \in \mathrm{jumpSet}(T) \iff
    \exists x,\; x \in T.\mathrm{level}(\tilde{n}) \;\land\; x \notin T.\mathrm{level}(\widetilde{n+1})
  \]
\end{proposition}

\begin{proof}
  定義の展開。\texttt{Set.Nonempty}，\texttt{Set.mem\_diff} を展開すれば等価。
\end{proof}

\subsection{Hasse-Arf 条件}

\begin{definition}[Hasse-Arf 条件]\label{def:hasse-arf-condition}
  分岐データ $\mathrm{rd}$，サイズ関数 $\mathrm{gs}$，$|G_0|$ のサイズ $g_0$ に対し、
  \textbf{Hasse-Arf 条件}を
  \[
    \mathrm{HasseArfCondition}(\mathrm{rd},\mathrm{gs},g_0) :=
    \forall b \in \mathrm{jumpSet}(\mathrm{ramificationTower}\,\mathrm{rd}),\;
    g_0 \mid \Phi(b)
  \]
  で定める。
\end{definition}

\begin{mathinsight}
条件 $g_0 \mid \Phi(b)$（$g_0 = |G_0|$）は
\[
  \frac{\Phi(b)}{|G_0|} = \varphi(b) \in \mathbb{Z}
\]
と等価。すなわち「上付番号のジャンプが整数値」という古典的 Hasse-Arf 定理の
\textsc{StructureTower} 版定式化。

\textsc{StructureTower} 的意味：Herbrand \texttt{reindex} による
$\mathrm{gradedPiece}$ の「移動先」が $\mathbb{N}$ の格子点に着地する。
アーベル拡大でないと、着地点が有理数になり、
$\mathbb{N}$ 上の \textsc{StructureTower} として整合的に定義できない。
\end{mathinsight}

\begin{proposition}[ジャンプ 0 での自明性]\label{prop:hasse-arf-zero}
  $\Phi(0) = 0$ は任意の $g_0$ で割り切れる：$g_0 \mid \Phi(0)$。
\end{proposition}

\begin{proof}
  \texttt{scaledHerbrand\_zero} より $\Phi(0) = 0$。\texttt{dvd\_zero}。
\end{proof}

\subsection{有限ジャンプ集合}

\begin{definition}[有限ジャンプ集合]\label{def:finite-jump-set}
  \[
    \mathrm{HasFiniteJumpSet}(\mathrm{rd}) :=
    \exists N \in \mathbb{N},\; \forall n \geq N,\; n \notin \mathrm{jumpSet}(\mathrm{ramificationTower}\,\mathrm{rd})
  \]
\end{definition}

\begin{theorem}[eventual bot から有限ジャンプ]\label{thm:finite-jump-eventual-bot}
  $\exists N,\;\forall n \geq N,\;G_n = \{1\}$ ならば $\mathrm{HasFiniteJumpSet}(\mathrm{rd})$。
\end{theorem}

\begin{proof}
  $N$ 以上では $G_n = \{1\}$，$G_{n+1} = \{1\}$ なので
  $\mathrm{gradedPiece}(T,n) = \{1\} \setminus \{1\} = \emptyset$。
  したがって $n \notin \mathrm{jumpSet}(T)$。

  形式的証明のキー：
  \begin{enumerate}
    \item $n \in \mathrm{jumpSet}(T)$ を仮定し矛盾を導く。
    \item 成員 $x$ を得て、$G_n = \{1\}$ より $x = 1$。
    \item $G_{n+1} = \{1\}$ より $1 \in G_{n+1}$，$x \notin G_{n+1}$ に矛盾。
  \end{enumerate}
\end{proof}

\begin{figure}[H]
\centering
\begin{tikzpicture}[>=stealth]
  \draw[->] (0,0) -- (9,0) node[right] {$n$};
  % Jumping region
  \fill[blue!15] (0,0.1) rectangle (5,1.5);
  % Stable region
  \fill[gray!20] (5,0.1) rectangle (8.5,0.4);
  \draw[very thick, red] (5, 0) -- (5, 1.7) node[above] {\small $N$};
  % Graded pieces in jumping region
  \foreach \x/\h in {0.5/1.4, 1.5/1.0, 2.5/0.7, 3.5/0.4, 4.5/0.2} {
    \fill[red!40] (\x-0.4, 0.1) rectangle (\x+0.4, \h);
  }
  \node[below, blue] at (2.5, -0.2) {\small ジャンプあり};
  \node[below, gray] at (6.5, -0.2) {\small ジャンプなし ($G_n = \{1\}$)};
  % Brace
  \draw[decorate, decoration={brace, amplitude=5pt}]
    (0, 1.8) -- (5, 1.8) node[midway, above=5pt] {\small 有限個のジャンプ};
\end{tikzpicture}
\caption{定理~\ref{thm:finite-jump-eventual-bot}：分離条件で安定すると、
十分大きい $n$ ではジャンプが消える。ジャンプ集合は有限。}
\label{fig:finite-jumps}
\end{figure}

% =============================================================================
\section{全体のまとめと展望}
% =============================================================================

\subsection{NT2 で構築したこと}

\begin{enumerate}[label=\S NT2-\arabic*.]
  \item \textbf{次数付き塔（Graded Tower）}：
    $\mathrm{gradedPiece}(T,n) = G_n \setminus G_{n+1}$。
    互いに素性（定理~\ref{thm:graded-piece-disjoint}），
    和集合との関係（命題~\ref{prop:union-graded}），
    分岐塔での具体化（命題~\ref{prop:graded-ramification}）。
    $\to$ \textsc{OrderHom} では自然に記述できない「隣接レベルの概念」。

  \item \textbf{Herbrand 関数}：
    $\mathrm{scaledHerbrand}\,\mathrm{gs}\colon \mathbb{N} \to \mathbb{N}$
    （$\Phi(n) = \sum \mathrm{gs}(i+1)$）。
    単調性（定理~\ref{thm:scaled-herbrand-mono}），
    狭義単調性（定理~\ref{thm:scaled-herbrand-strict-mono}），
    Finset.sum 表現，
    $\mathrm{herbrandReindex}\colon \mathbb{N}^{\mathrm{op}} \to \mathbb{N}^{\mathrm{op}}$。

  \item \textbf{上付番号塔}：
    $\mathrm{upperNumberingTower}(\mathrm{rd},\mathrm{gs})
    = \mathrm{reindex}\,(\Phi)\,(\mathrm{ramificationTower}\,\mathrm{rd})$。
    L1 の \texttt{reindex} の\textbf{最初の数論的インスタンス}。
    レベル 0 の整合性（命題~\ref{prop:upper-zero}），
    NatInclusion 関係（定理~\ref{thm:upper-nat-inclusion}），
    二重 reindex 合成則（定理~\ref{thm:double-reindex}）。

  \item \textbf{商との整合性}：
    $\mathrm{restrictToSubgroup}$：塔の部分群への制限。
    $\mathrm{UpperNumberingCompatible}$：上付番号が制限と可換である条件の定式化。
    $\mathrm{lowerNumbering\_is\_trivial\_reindex}$（定理~\ref{thm:lower-is-id-reindex}）。
    $\to$ \textbf{Tower Comparison Theory の数論的基本定理}。

  \item \textbf{Hasse-Arf の定式化}：
    $\mathrm{jumpSet}(T)$：塔の段差が存在するインデックス集合。
    $\mathrm{HasseArfCondition}$：$|G_0| \mid \Phi(b)$（全ジャンプに対して）。
    $\mathrm{HasFiniteJumpSet}$：分岐塔のジャンプ集合は有限。
    $\to$ アーベル拡大でのみ上付番号が $\mathbb{N}$ 上で整合する。
\end{enumerate}

\subsection{StructureTower 理論への貢献（NT1 + NT2）}

\begin{table}[H]
\centering
\caption{NT1 + NT2 の貢献と \textsc{OrderHom} との比較}
\label{tab:contributions}
\renewcommand{\arraystretch}{1.3}
\begin{tabular}{p{3.5cm}p{2.5cm}p{5.5cm}}
\toprule
\textbf{発見} & \textbf{OrderHom で可能か} & \textbf{StructureTower の優位性} \\
\midrule
gradedPiece & \mbox{$\triangle$（定義は可能）} & 隣接概念が添字の離散構造に依存 \\
jumpSet & $\triangle$ & gradedPiece から自然に導出 \\
Herbrand reindex & $\bigcirc$（関数合成） & 「塔の内部構造が reindex を決定」 \\
商との整合性 & $\times$ & restrictToSubgroup + NatInclusion \\
二重 reindex & $\bigcirc$ & reindex\_comp の直接適用 \\
Hasse-Arf 条件 & $\times$ & jumpSet + scaledHerbrand の整除性 \\
\bottomrule
\end{tabular}
\end{table}

特に「商との整合性」（\S NT2-4）は \textsc{OrderHom} の語彙では
自然に記述できない。「異なる対象上の2つの \textsc{OrderHom} を
共通の添字空間で比較し，部分構造への制限と可換になる reindex を見つける」
という問題は，\textsc{StructureTower} の $\mathrm{Hom} + \mathrm{reindex} + \mathrm{NatInclusion}$
の組み合わせで初めて自然に定式化できる。

\subsection{NT3 以降の展望}

\begin{figure}[H]
\centering
\begin{tikzpicture}
  \node[draw,rounded corners,fill=gray!10,text width=3cm,align=center] (nt1)
    {\textbf{NT1}\\分岐塔\\下付番号};
  \node[draw,rounded corners,fill=blue!10,text width=3.5cm,align=center,
    right=1.2cm of nt1] (nt2)
    {\textbf{NT2}\\Herbrand\\上付番号\\（本稿）};
  \node[draw,rounded corners,fill=green!10,text width=4cm,align=center,
    right=1.2cm of nt2] (nt3)
    {\textbf{NT3}\\局所類体論\\Artin 写像\\Tower Comparison};
  \node[draw,rounded corners,fill=orange!10,text width=3.5cm,align=center,
    below=1.2cm of nt2] (nt4)
    {\textbf{NT4}\\副有限群\\$\pi_1(X)$ の塔\\遠アーベル幾何};
  \draw[->] (nt1) -- (nt2);
  \draw[->] (nt2) -- (nt3);
  \draw[->,dashed] (nt3) -- (nt4);
  \draw[->,dashed] (nt2) -- (nt4);
\end{tikzpicture}
\caption{数論シリーズの全体展望。NT2（本稿）は下付番号の欠陥を修正し，
類体論（NT3）と遠アーベル幾何（NT4）への橋を準備する。}
\label{fig:nt-roadmap}
\end{figure}

\begin{description}
  \item[NT3: 局所類体論と Tower Comparison]
    Artin 写像 $\theta\colon K^\times \to G^{\mathrm{ab}}$ を通じて，
    \texttt{idealPowTower}（$K$ の素イデアル冪）と
    \texttt{upperNumberingTower}（$G$ の上付番号分岐）を
    $\theta$ が \textsc{StructureTower} Hom として結ぶことを形式化。
    $\to$ Tower Comparison Theory の完全な数論的インスタンス。

  \item[NT4: 副有限群と逆極限]
    $\pi_1(X)$ のフィルトレーション $\to$ \textsc{StructureTower}。
    逆極限 $= \mathrm{global}$ の副有限版。
    望月の「宇宙際 Teichmüller 理論」= 異なる副有限塔間の「リンク」？
\end{description}

% =============================================================================
\end{document}
